\pg{357} \emph{a) O ente, que é acto puro, é immutavel.} --- O ente, que é \emph{acto puro}, não contem em si nenhuma potencialidade. Ora o ente, que não contem em si nenhuma potencialidade, é immutavel; porque toda a mudança importa uma passagem da potencia para o acto. Logo o ente, que é \emph{acto puro}, é \emph{immutavel}.

\emph{b) O ente, que é immutavel, é acto puro.} --- O ente, que é \emph{immutavel}, não contem em si nenhuma potencialidade. Ora o ente, que não contem em si nenhuma potencialidade, é acto puro. Logo o ente, que é \emph{immutavel}, é \emph{acto puro} (¹).

\textbf{86. O ente mutavel não pode ser acto puro.} --- O ente \emph{mutavel} em tanto pode estar sujeito a mudanças, em quanto contem em si alguma potencialidade. Ora o que contem em si alguma potencialidade não é acto puro. Logo o ente \emph{mutavel} não pode ser \emph{acto puro}.\footnote{O ente que se muda, muda-se emquanto perde alguma coisa e adquire outra. Ora esta perda e esta acquisição não seria possivel, se não se desse uma \emph{composição} do sujeito que se muda, --- \emph{antes} da mudança, com a coisa que se perde, --- \emph{depois} da mudança, com a coisa que se adquire. Diz S. Thomaz: «Omne quod movetur, quantum ad aliquid manet, et quantum ad aliquid transit; sicut quod movetur de albedine in nigredinem, manet secundum substantiam. Et sic in omni quod movetur, attenditur aliqua compositio» (\emph{Sum}. \emph{Th}., q. I, q. 9, a. 1). Se \emph{o} ente mutavel é necessariamente \emph{composto}, não é \emph{acto puro}.}

\textbf{87. O ente immutavel é necessario, infinito, simplicissimo; o ente necessario, infinito, simplicissimo, é immutavel.}

\emph{a) O ente immutavel é necessario, infinito, simplicissimo.} --- O ente immutavel é acto puro. Ora o ente, que é acto puro, é necessario, infinito, simplicissimo. Logo o ente \emph{immutavel} é \emph{necessario}, \emph{infinito}, \emph{simplicissimo}.

\emph{b) O ente necessario, infinito, simplicissimo, é immutavel.} ---\footnote{Ninguem se admire se a demonstração d' estas proposições se basêa sempre neste criterio --- se um ente tem, ou não, \emph{potencialidade}. A essencia de tudo o que existe ou é \emph{acto} ou é composta de \emph{potencia} e de \emph{acto}. E, como a essencia é a raiz de todas as propriedades, é claro que a ella é preciso recorrer, como ao verdadeiro e intrinseco criterio, para conhecermos o que convem a um ente. --- \emph{O} ente, que é \emph{acto puro}, exclue toda a potencialidade, e por isso não é capaz nem de perder nem de adquirir alguma coisa, é \emph{immutavel}. --- O ente, que é \emph{immutavel}, não é capaz de adquirir ou de perder coisa alguma, por isso não tem potencialidade, é \emph{acto puro}. --- Diz S. Thomaz: «Primum ens oportet esse purum actum absque permixtione alicujus potentiae, eo quod potentia simpliciter est posterior actu. Omne autem quod quocumque modo mutatur, est aliquo modo in potentia. Ex quo patet quod impossibile est \textsc{Deum} aliquo modo mutari» (\emph{Sum}. \emph{Th}., q. I, q. 9, a. 1).}